\clearpage
\VerdierPhysicalPageAnchor{0019}

\VerdierSemanticAnchor{intro-1-resolution-comparison-diagram}
\[
\begin{tikzcd}[column sep=large,row sep=large]
X \arrow[r,"\varepsilon'"] \arrow[dr,"\varepsilon"']
  & \mathcal R(X) \arrow[d,"\rho"] \\
  & I^{\bullet}(X)
\end{tikzcd}
\]

\VerdierSemanticAnchor{intro-1-f-acyclic-morphism}
さらに、$F$-非輪状分解の射

\VerdierSemanticAnchor{intro-1-f-acyclic-morphism-diagram}
\[
\begin{tikzcd}[column sep=large,row sep=large]
X \arrow[r,"\varepsilon_1"] \arrow[dr,"\varepsilon_2"']
  & \mathcal R_1(X) \arrow[d,"\rho'"] \\
  & \mathcal R_2(X)
\end{tikzcd}
\qquad\text{,}
\]
が与えられると、$\mathcal B$ の対象の複体の射

\VerdierSemanticAnchor{intro-1-f-rho-prime}
\[
F(\rho') : F(\mathcal R_1(X)) \longrightarrow F(\mathcal R_2(X))
\]
は擬同型である。

\VerdierSemanticAnchor{intro-1-inductive-system}
ここでは、$\mathcal B$ において複体と擬同型からなる帰納系が得られ、
それはホモトピーを除いて一意な最大元、すなわち $X$ の入射分解の $F$ に
よる像をもつことに注意しよう。これから扱う第二の例では、この現象は
もはや生じない。

\VerdierSemanticAnchor{intro-1-left-derived-tensor}
分解によって導来関手を定義する方法は、$\mathcal A$ が十分な入射対象や
射影対象をもつとは限らない場合にも、ときに一般化できる。$\mathcal A$ を
位相空間 $E$ 上のアーベル群の層の圏としよう。この圏は一般には十分な
射影対象をもたない。それでも、各層 $G$ に層
$F\otimes_{\mathbb Z}G$（$F$ は固定した層）を対応させる
$\mathcal A$ から $\mathcal A$ への右完全関手の\emph{左導来関手}を
定義したい。当然、この定義は左導来関手の通常の性質（完全列に対する
振る舞いと、ある普遍性）をもつ関手を与え、既知の場合（たとえば $E$ が
一点に縮退する場合）には通常の Tor 積を再現しなければならない。

\VerdierSemanticAnchor{intro-1-left-resolution}
次のようにすればよい。各層 $G$ は有限長（実際には長さ 2）の左分解
$L^{\bullet}\to G$ をもち、その各層の
